\subsection{Distribución de clases}
Hay 5174 clientes con \var{Churn=No} (73.46\%) y 1869 con \var{Churn=Yes} (26.54\%), como muestra 
la Fig.~\ref{fig:balance}. La clase positiva es minoritaria, pero ese hecho no determina por sí 
solo cuán difícil será clasificarla. Un modelo que siempre prediga No obtendría 73.46\% 
de exactitud, pero no detectaría ningún abandono. Por eso, la evaluación también debe considerar 
cuántos abandonos detecta el modelo y cuántas falsas alertas genera.

\begin{figure}[t]
    \centering\includegraphics[width=\columnwidth]{balance.png}
    \caption{Distribución del abandono en los 7043 clientes. Las etiquetas muestran cantidad y 
    porcentaje de cada clase.}
    \label{fig:balance}
\end{figure}

\subsection{Distribuciones numéricas y comparación por abandono}
La Tabla~\ref{tab:numericas} presenta las estadísticas descriptivas. La antigüedad mediana es 29 meses 
y su mitad central se encuentra entre 9 y 55 meses. En el cargo acumulado, la media de 2283.30 supera a 
la mediana de 1397.48, en concordancia con una distribución extendida hacia montos altos. El cargo mensual 
presenta el orden contrario, con media de 64.76 y mediana de 70.35.

\begin{table}[t]
    \caption{Resumen de variables cuantitativas}
    \label{tab:numericas}\centering\footnotesize
    \begin{tabular}{@{}lrrr@{}}\toprule
        Medida & Antigüedad & Cargo mensual & Cargo acumulado\\\midrule
        Observaciones & 7043 & 7043 & 7032\\
        Media & 32.37 & 64.76 & 2283.30\\
        Mediana & 29.00 & 70.35 & 1397.48\\
        Desv. estándar & 24.56 & 30.09 & 2266.77\\
        Q1 & 9.00 & 35.50 & 401.45\\
        Q3 & 55.00 & 89.85 & 3794.74\\
        Mínimo & 0.00 & 18.25 & 18.80\\
        Máximo & 72.00 & 118.75 & 8684.80\\\bottomrule
    \end{tabular}
\end{table}

Entre quienes abandonaron, la antigüedad mediana es 10 meses, frente a 38 meses entre quienes no 
abandonaron. Su cargo mensual mediano es mayor (79.65 frente a 64.43), pero su cargo acumulado mediano 
es menor (703.55 frente a 1683.60). Una persona puede pagar más por mes y acumular menos si lleva menos 
tiempo contratando el servicio. Las cajas de la Fig.~\ref{fig:boxplots} también se superponen, por lo que 
estas características no separan perfectamente a quienes abandonaron de quienes no lo hicieron.

\begin{figure*}[t]
    \centering\includegraphics[width=\textwidth]{boxplots.png}
    \caption{Antigüedad, cargo mensual y cargo acumulado según abandono. No y Yes indican ausencia y 
    presencia de abandono registrado. Los puntos fuera de los bigotes se determinan dentro de cada grupo.}
    \label{fig:boxplots}
\end{figure*}

Los histogramas muestran concentraciones de quienes abandonaron en antigüedades y cargos 
acumulados bajos. Muestran cantidades de clientes, así que las diferencias de altura también reflejan el 
distinto tamaño de los grupos. Por eso, no deben interpretarse como probabilidades de abandono por intervalo.

\subsection{Correlaciones entre variables cuantitativas}

Como se observa en la Fig.~\ref{fig:pearson}, la correlación de Pearson entre antigüedad y 
cargo acumulado es 0.826; entre cargo mensual y acumulado, 0.651; y entre antigüedad y 
cargo mensual, 0.248. El acumulado está asociado con el tiempo y los pagos, pero no
necesariamente equivale a multiplicar los meses por el cargo mensual actual. Estas asociaciones 
no justifican eliminar automáticamente una variable ni constituyen por sí solas una fuga 
de información.

Para cada correlación se utilizan los clientes con datos disponibles en ambas variables. 
Las comparaciones que incluyen el cargo acumulado consideran 7032 clientes, mientras que 
la comparación entre antigüedad y cargo mensual utiliza los 7043.

\begin{figure}[t]
    \centering
    \includegraphics[width=\columnwidth]{pearson.png}
    \caption{Matriz de correlación de Pearson entre antigüedad,
    cargo mensual y cargo acumulado.}
    \label{fig:pearson}
\end{figure}
